% ======================================================================
% sections/discussion.tex
%
% Discussion scaffold for the paper:
%   Triple Humanoid 24/7 Adverse Event Oncology Trial Response Team:
%   4-Site Rotation
%
% This file is a draft scaffold. A future Claude Code Opus 4.7 1M Max
% session reads the bracketed instruction blocks below, then reads the
% files those instructions name, then expands the scaffold into the
% finished discussion. Target length is roughly 10 to 14 pages of the
% 70 plus page final paper.
% ======================================================================

\begin{sectionaccent}

[DISCUSSION INSTRUCTIONS FOR FUTURE CLAUDE CODE OPUS 4.7 1M MAX SESSION.
Write the discussion as 6 subsections inside this single
\texttt{\textbackslash section\{Discussion\}} block. Tie every claim back to a
results subsection.

Subsection 1, Thesis Reconciliation. Restate the thesis verbatim from
\texttt{demo-projects/07-humanoid/paper/instructions/README.md} and
walk through how each clause of the thesis is supported by:
\begin{itemize}
  \item ``On-premises repository based LLMs provide commands to
        humanoid robots'' is supported by
        \texttt{codegen/src/llm\_planner.py},
        \texttt{codegen/src/broadcaster.py}, and
        \texttt{codegen/src/comms\_intellectual.py}.
  \item ``based on real-time sensor data'' is supported by
        \texttt{codegen/src/sensor\_to\_xyz.py} and
        \texttt{codegen/src/peer\_state\_tracker.py}.
  \item ``controlled via x, y, z coordinates'' is supported by
        \texttt{codegen/src/cartesian\_planner.py},
        \texttt{codegen/src/kinematics.yaml}, and
        \texttt{codegen/src/xyz\_safety.py}.
  \item ``to administer synergistic treatment to patient adverse
        events'' is supported by
        \texttt{codegen/src/ctcae\_grader.py} and
        \texttt{codegen/src/physician\_escalation.py}.
  \item ``This workflow minimizes single robot error potential'' is
        supported by the 191 line
        \texttt{codegen/src/swarm\_coordinator.py} (the camaraderie
        factor of 3).
\end{itemize}

Subsection 2, Significance Of The On-Premises Repository LLM Pattern.
Discuss the quality and safety benefits of having one Claude Opus 4.7
1M instance per site rather than a remote cloud. Quality: every
inference traces to a known on-prem repository, which simplifies the
audit trail for the FDA RTCT pilot. Safety: a network partition
between the site and the public internet does not stop the 1 Hz
broadcast cadence. Speed: the 200 ms median latency named in the
instructions README is achievable on a local LAN. Read:
\begin{itemize}
  \item \texttt{demo-projects/07-humanoid/paper/instructions/README.md}
        Inventory at a Glance table.
  \item \texttt{demo-projects/07-humanoid/paper/codegen/src/llm\_planner.py}.
  \item \texttt{demo-projects/07-humanoid/paper/codegen/src/prompt\_frozen.md}.
\end{itemize}

Subsection 3, Significance Of The Cartesian Coordinate Surface. The
LLM does not author joint angles; it authors x, y, z goals and lets
the kinematics layer convert. This is the contract that lets a
regulator review the safety envelope from the constants alone. Read:
\begin{itemize}
  \item \texttt{demo-projects/07-humanoid/paper/codegen/src/kinematics.yaml}.
  \item \texttt{demo-projects/07-humanoid/paper/codegen/src/cartesian\_planner.py}.
  \item \texttt{demo-projects/07-humanoid/paper/codegen/src/xyz\_safety.py}.
  \item \texttt{demo-projects/07-humanoid/paper/codegen/tests/test\_xyz\_safety.py}.
\end{itemize}

Subsection 4, Significance Of The 32 Iteration Sweep And The
Competition Framing. Read:
\begin{itemize}
  \item \texttt{demo-projects/07-humanoid/paper/codegen/data/iterations/index.jsonl}
        (32 iterations across 4 sites and many columns).
  \item \texttt{demo-projects/07-humanoid/paper/codegen/src/iterate.py}
        and \texttt{demo-projects/07-humanoid/paper/codegen/src/compare\_agent.py}.
  \item Author's December 2025 paper
        \texttt{kawchak\_2025\_18029100}: ``Code Generation
        Competition: 16 Proprietary vs. Open-Source LLMs and Iterative
        Learning Based on FDA Adverse Event Reporting System.''
\end{itemize}
Tie the 32 iteration sweep back to the prior 16 LLM code generation
competition. The earlier competition established LLM effectiveness
against FDA adverse event data. The current 32 iteration sweep
extends that evidence to robot dispatch and Cartesian command code.

Subsection 5, FDA, CTCAE, And Other Regulators. Read:
\begin{itemize}
  \item \texttt{demo-projects/07-humanoid/paper/codegen/src/fda\_rtct\_submitter.py}
        (67 lines, HR 9505 SLA, ed25519 signing in production,
        placeholder in simulation).
  \item \texttt{demo-projects/07-humanoid/paper/codegen/src/ctcae\_grader.py}.
  \item \texttt{demo-projects/07-humanoid/paper/codegen/src/check\_errors.py}
        (LLM credibility in error checking code).
  \item \texttt{demo-projects/07-humanoid/paper/codegen/tests/test\_xyz\_safety.py}
        (LLM credibility in test code).
\end{itemize}
Write this subsection respectfully but advocate for accelerating
review timelines. The United States needs to remain Number 1 in the
world regarding patient safety, efficacy, and speed benefits in
oncological adverse event robot swarms for clinical trials. Frame the
review work as something the project has already checked off so the
agencies can focus on policy and not on rebuilding the audit trail.

Subsection 6, Prior Paper Context Integration. Read both Deep-Research
trees in
\texttt{demo-projects/07-humanoid/paper/inputs/Deep-Research-A/} and
\texttt{demo-projects/07-humanoid/paper/inputs/Deep-Research-B/}:
\begin{itemize}
  \item Deep-Research-A is the Medical Full Automation A source
        bundle: \texttt{README.md},
        \texttt{chunk1\_paper\_text.md},
        \texttt{chunk2\_tables\_references.md},
        \texttt{chunk3\_bibtex.md}.
  \item Deep-Research-B is the Medical Rapid Prototyping B source
        bundle: \texttt{chunk\_01\_executive\_summary\_and\_fully\_automated\_sponsor.md},
        \texttt{chunk\_02\_federated\_learning\_and\_patient\_prediction.md},
        \texttt{chunk\_03\_clinagent.md},
        \texttt{chunk\_04\_comparative\_overview\_and\_implications.md},
        \texttt{chunk\_05\_bibtex\_references.md}.
\end{itemize}
Compare and contrast the prior paper's Medical Full Automation A and
Medical Rapid Prototyping B framings with the current paper's 4 site
rotation framing. Position the current paper as the natural successor.

Required real-life feasibility insight to include in this section:
``The on-premises repository LLM pattern is feasible today on a single
high end server per site. The harder gap to close is not compute; it
is the IRB, sponsor, and site coordination required to enroll the
first real PAT-NET-001 cohort. The paper is a checklist of the
technical work that has been completed so that the human work of
governance can move forward without waiting on the technology.''

Do not exceed 14 pages. Use single dashes only. Use \S\ for the
section sign. Avoid stranding 1 or 2 word lines on their own.]

The discussion section will reconcile each clause of the thesis with
specific files, walk through the quality and safety significance of
the on-premises repository LLM pattern and the Cartesian coordinate
surface, position the 32 iteration sweep against the prior 16 LLM
competition, address FDA CTCAE and other regulators respectfully, and
integrate the prior paper context from Deep-Research-A and
Deep-Research-B.

\sectiontable{Discussion planning matrix}{%
Thesis & D1 & 1 & instructions README & Per clause source list & 2 pages & Every clause sourced & Ready \\
OnPrem & D2 & 2 & llm\_planner plus prompt & Quality, safety, speed & 2 pages & Audit trail framed & Ready \\
Cartesian & D3 & 3 & kinematics plus planner & Coordinate contract & 2 pages & Constants only review & Ready \\
Iteration & D4 & 4 & iterations index & 32 sweep significance & 2 pages & Cites kawchak\_2025 & Ready \\
Regulator & D5 & 5 & fda\_rtct plus ctcae & Respectful advocacy & 2 pages & HR 9505 named & Ready \\
PriorPaper & D6 & 6 & Deep-Research A and B & Successor framing & 2 to 3 pages & Both bundles named & Ready \\
}

\end{sectionaccent}
